\section*{Materials and methods}

\subsection*{Training collection and labels}
The training collection is the Arcadia Science \Ecoli{} genotype--phenotype release
\cite{arcadia2024dataset,bell2025arcadia}, distributed in two parts: variant calls against a
pan-genome reference with a phenotype table and strain metadata, and a gene presence/absence
release in PLINK format \cite{purcell2007plink}. A strain was retained when its sequencing
accession appeared in the variant file, the mapping file gave a genome name, and that name
appeared in the phenotype table, leaving 6,962 strains used in every analysis. Strain metadata
matched for 6,821 of them (98.0\%).

Of the sixteen antibiotics in the phenotype table we kept the seven with workable label counts and
class balance: ampicillin (AMP), ciprofloxacin (CIP), amoxicillin--clavulanic acid (AMC),
trimethoprim--sulfamethoxazole (SXT), gentamicin (GEN), cefotaxime (CTX) and ceftazidime (CAZ),
spanning five drug classes. Phenotypes were coded 1 for resistant and 0 for susceptible, with
intermediate results counted as resistant, following the authoring group and the clinical reading
that an intermediate result makes the standard dose insufficient; intermediates were excluded in a
sensitivity analysis. No label was imputed: a $6{,}962 \times 7$ mask records which entries were
measured, and every loss and metric is computed over observed entries only. Coverage ranges from
49.1\% (SXT) to 86.7\% (CIP, GEN) and resistance prevalence from 0.118 (GEN) to 0.501 (AMP)
(Table~\ref{tab:labels}). Both spreads matter downstream: fewer labels widen intervals, and because
the no-skill value of average precision is the prevalence, a raw score cannot be compared across
drugs.

Collapsing intermediate into resistant removes the three-category comparison, so categorical
agreement reduces to binary accuracy throughout, and our error rates are not the quantities a
susceptibility-testing standard tabulates.

\subsection*{Feature representations}
\textbf{SNP blocks by gene.} The compressed variant file (2,451,184 variants, matching the figure
the authoring group reports) was streamed once; indels were dropped (43,799) and variants with a
minor allele frequency below 1\% removed (1,871,946), leaving 535,439. Genotypes are kept as small
integers (0 reference, 1--3 alternate; 83.3\% of calls are reference and 2.5\% of non-reference
calls are second or third alternate, too few to justify one-hot expansion). Missing calls were set
to the reference allele, as in the source release. This is not neutral: coverage is lowest on
plasmids and other mobile elements that assemble poorly from short reads, which is where acquired
resistance genes sit, and the Discussion returns to it. Two gene-level filters, both fitted on the
training partition of whichever split is being evaluated, reduce the matrix: a variability filter
keeping genes with at least five SNPs (18,975 $\rightarrow$ 13,799 on the random split) and an
importance filter keeping the 5,000 genes with the highest summed gradient-boosting gain over the
seven per-drug models (each trained on that drug's top 50,000 $\chi^2$-filtered SNPs). The 5,000
retained genes carry 262,225 SNPs, 49.0\% of the filtered variants. Genes are read as fixed blocks
of 64 positions (median 32 SNPs per gene, mean 52, maximum 717); 28\% of genes are truncated and
the rest zero-padded and masked.

The variant-level filters above were computed once over the whole collection before any split
existed. They use no phenotype label, so they cannot transmit outcome information, but the allele
frequencies of test strains do influence which columns exist. We state this rather than repeat the
filters per split, which would give the three splits different feature spaces and make the
comparison the paper is built on incoherent.

\textbf{Gene presence/absence.} The release encodes 32,441 pan-genome loci as 0 and 2 without
stating which value means present. The two readings are biologically opposite, so we resolved them
against an external check: under one reading the median locus prevalence is 0.966 and the average
strain carries about 28,700 genes; under the other the median is 0.034 and the average strain
carries 3,709 (range 2,141--8,333). An \Ecoli{} genome holds roughly 4,000--5,000 genes, which
selects the second reading. Rebuilding the matrix in that orientation flipped 10,511 of 32,441
loci and left 21,930 unchanged; 2,924 loci are monomorphic. After the rebuild the core genes
\textit{gyrA}, \textit{parC}, \textit{rpoB} and \textit{recA} sit at prevalence 1.0, and locus
prevalence correlates with the ECOR reference frequencies at $\rho = 0.689$. The wrong reading
would have raised no error, only inverted every gene, so a per-strain gene count is now checked
automatically on each matrix build. Of the 6,962 strains, 6,949 appear in the PLINK sample list;
the remaining 13 (0.19\%) are carried as all-zero vectors, which conflates a missing record with an
empty accessory genome. We record the defect rather than report a corrected set of results, as its
effect is below the resolution of our intervals.

\textbf{Per-gene statistics.} Each of the 5,000 genes is also summarised by four values computed
from its SNPs: allele mean, allele variance, number of non-reference positions and allele maximum,
standardised with training-partition means and standard deviations only.

\textbf{Lineage encoding.} Phylogroup, in the Clermont sense \cite{clermont2013phylotyping}, was
taken from the typing field of the strain metadata where present and otherwise from a look-up table
of 57 common sequence types. That table's provenance is not recorded in the preparation notebooks,
so we report it as an undocumented lookup and use the phylogroup only as a coarse covariate and
grouping variable; calling it from each genome with ClermonTyping \cite{beghain2018clermontyping}
is the established route and is what a study resting on the phylogroup should do. The composition
is dominated by B2 (50.9\%), with 24.4\% unassigned and phylogroup E holding four strains, so no
per-group estimate is attempted.

\textbf{Determinant calls.} For each strain we used the AMRFinderPlus
\cite{feldgarden2019amrfinder,feldgarden2021amrfinderplus} calls published by NCBI Pathogen
Detection \cite{ncbi_pathogen_detection} (software 4.2.7, database 2026-03-24.1), matched by
sequencing run and BioSample. Calls without a suffix were taken as acquired genes and
\texttt{=POINT} calls as point mutations; partial, end-of-contig, mistranslated and HMM-only calls
were not counted. Genes and mutations present in at least five strains were kept (89 acquired-gene
and 82 point-mutation columns). Calls were available for 5,072 of 6,962 strains (72.9\%); strains
without calls are not treated as determinant-free but are excluded from every analysis using calls.

\textbf{Metadata and assembly statistics.} Country, collection year, submitting study and
sequencing platform; and five assembly statistics: assembly length, coding-sequence count, GC
content, contig count and contig N50.

\subsection*{Feature integration checks}
Two checks run before any training. The first confirms that all 6,962 strain identifiers appear in
the same order in every view, since otherwise one genome's SNPs would be paired with another's
genes. The second compares gene identifiers across views, and it caught a real fault: the PLINK
reader had invented placeholder names instead of reading the annotation file, so the two views
shared no gene name at all. After the fix, all 5,000 genes selected for the SNP views are present
in the 32,441-gene catalogue. Neither fault would have produced a visible error.
